%%%%%%%%%%%%%%%%%%%%%%%%%%%%% Define Article %%%%%%%%%%%%%%%%%%%%%%%%%%%%%%%%%%
\documentclass{article}
%%%%%%%%%%%%%%%%%%%%%%%%%%%%%%%%%%%%%%%%%%%%%%%%%%%%%%%%%%%%%%%%%%%%%%%%%%%%%%%

%%%%%%%%%%%%%%%%%%%%%%%%%%%%% Using Packages %%%%%%%%%%%%%%%%%%%%%%%%%%%%%%%%%%
\usepackage[T1]{fontenc}
\usepackage[polish]{babel}
\usepackage{geometry}
\usepackage{graphicx}
\usepackage{amssymb}
\usepackage{amsmath}
\usepackage{amsthm}
\usepackage{empheq}
\usepackage{mdframed}
\usepackage{booktabs}
\usepackage{lipsum}
\usepackage{graphicx}
\usepackage{color}
\usepackage{psfrag}
\usepackage{pgfplots}
\usepackage{bm}
%%%%%%%%%%%%%%%%%%%%%%%%%%%%%%%%%%%%%%%%%%%%%%%%%%%%%%%%%%%%%%%%%%%%%%%%%%%%%%%

% Other Settings

%%%%%%%%%%%%%%%%%%%%%%%%%% Page Setting %%%%%%%%%%%%%%%%%%%%%%%%%%%%%%%%%%%%%%%
\geometry{a4paper}

%%%%%%%%%%%%%%%%%%%%%%%%%% Define some useful colors %%%%%%%%%%%%%%%%%%%%%%%%%%
\definecolor{ocre}{RGB}{243,102,25}
\definecolor{mygray}{RGB}{243,243,244}
\definecolor{deepGreen}{RGB}{26,111,0}
\definecolor{shallowGreen}{RGB}{235,255,255}
\definecolor{deepBlue}{RGB}{61,124,222}
\definecolor{shallowBlue}{RGB}{235,249,255}
%%%%%%%%%%%%%%%%%%%%%%%%%%%%%%%%%%%%%%%%%%%%%%%%%%%%%%%%%%%%%%%%%%%%%%%%%%%%%%%

%%%%%%%%%%%%%%%%%%%%%%%%%% Define an orangebox command %%%%%%%%%%%%%%%%%%%%%%%%
\newcommand\orangebox[1]{\fcolorbox{ocre}{mygray}{\hspace{1em}#1\hspace{1em}}}
%%%%%%%%%%%%%%%%%%%%%%%%%%%%%%%%%%%%%%%%%%%%%%%%%%%%%%%%%%%%%%%%%%%%%%%%%%%%%%%

%%%%%%%%%%%%%%%%%%%%%%%%%%%% English Environments %%%%%%%%%%%%%%%%%%%%%%%%%%%%%
\newtheoremstyle{mytheoremstyle}{3pt}{3pt}{\normalfont}{0cm}{\rmfamily\bfseries}{}{1em}{{\color{black}\thmname{#1}~\thmnumber{#2}}\thmnote{\,--\,#3}}
\newtheoremstyle{myproblemstyle}{3pt}{3pt}{\normalfont}{0cm}{\rmfamily\bfseries}{}{1em}{{\color{black}\thmname{#1}~\thmnumber{#2}}\thmnote{\,--\,#3}}
\theoremstyle{mytheoremstyle}
\newmdtheoremenv[linewidth=1pt,backgroundcolor=shallowGreen,linecolor=deepGreen,leftmargin=0pt,innerleftmargin=20pt,innerrightmargin=20pt,]{theorem}{Theorem}[section]
\theoremstyle{mytheoremstyle}
\newmdtheoremenv[linewidth=1pt,backgroundcolor=shallowBlue,linecolor=deepBlue,leftmargin=0pt,innerleftmargin=20pt,innerrightmargin=20pt,]{definition}{Definition}[section]
\theoremstyle{myproblemstyle}
\newmdtheoremenv[linecolor=black,leftmargin=0pt,innerleftmargin=10pt,innerrightmargin=10pt,]{problem}{Problem}[section]
%%%%%%%%%%%%%%%%%%%%%%%%%%%%%%%%%%%%%%%%%%%%%%%%%%%%%%%%%%%%%%%%%%%%%%%%%%%%%%%

%%%%%%%%%%%%%%%%%%%%%%%%%%%%%%% Plotting Settings %%%%%%%%%%%%%%%%%%%%%%%%%%%%%
\usepgfplotslibrary{colorbrewer}
\pgfplotsset{width=8cm,compat=1.9}
%%%%%%%%%%%%%%%%%%%%%%%%%%%%%%%%%%%%%%%%%%%%%%%%%%%%%%%%%%%%%%%%%%%%%%%%%%%%%%%

%%%%%%%%%%%%%%%%%%%%%%%%%%%%%%% Title & Author %%%%%%%%%%%%%%%%%%%%%%%%%%%%%%%%
\title{Badanie efektywności algorytmów grafowych w zależności od rozmiaru grafu i sposobu reprezentacji}
\author{Krystian Mikołajczak 284452}
%%%%%%%%%%%%%%%%%%%%%%%%%%%%%%%%%%%%%%%%%%%%%%%%%%%%%%%%%%%%%%%%%%%%%%%%%%%%%%%

\begin{document}
    \maketitle
    \section{Wstęp}
    \subsection{Opis zadania}
      Celem było zaimplementowanie wybranych algorytmów rozwiązujących problemy grafowe. Rozwiązywane problemy to minimalne drzewo rozpinające i znalezienie najkrótszej ścieżki w grafie.
      Wyniki efektywności rozwiązywania zadanych problemów za pomocą wybranych algorytmów następnie zostały poddane analizie oraz porównane między sobą. Badana będzie efektywność algorytmów w zależności 
      od ilości wierzchołków w grafie i od gęstości krawędzi w grafie.
    \subsection{Algorytmy poddane analizie}
      Analizie poddano dwa algorytmy dla każdego z dwóch problemów. Każdy z algorytmów został zbadany dla reprezentacji grafu w formie macierzy sąsiedztwa oraz w formie listy następników/poprzedników.  
    \subsubsection{Minimalne drzewo rozpinające}
      Jest do drzewo zawierające wszystkie wierzchołki danego grafu o najmniejszej mozliwej sumie wag krawędzi.
      \begin{itemize}
      \item{Algorytm Prima}\\
        Algorytm Prima buduje strukturę krok po kroku, rozpoczynając od wybranego wierzchołka i zachłannie dołączając najlżejsze krawędzie łączące odwiedzoną część grafu z nieodwiedzoną. Implementacja wykorzystująca macierz sąsiedztwa wymaga przeglądania całego wiersza (wszystkich $V$ wierzchołków) przy każdym zdjęciu elementu z kolejki priorytetowej, co daje pesymistyczną złożoność rzędu $O(V^2 \log V)$. Zastosowanie listy następników pozwala z kolei na iterowanie wyłącznie po rzeczywistych krawędziach, co obniża czas do $O(E \log V)$. Zakłada się więc, że reprezentacja listowa będzie znacznie szybsza dla grafów rzadkich, podczas gdy macierzowa może być konkurencyjna w grafach gęstych z uwagi na mniejszy narzut pamięciowy i ciągły układ w pamięci.
        \cite{Prim}
      \item{Algorytm Kruskala}\\
        Algorytm Kruskala początkowo wyodrębnia wszystkie krawędzie grafu, sortuje je według wag w kolejce priorytetowej, a następnie łączy wierzchołki, korzystając ze struktury zbiorów rozłącznych (find-union) w celu unikania cykli. Sposób reprezentacji grafu determinuje narzut na początkowe pozyskanie krawędzi. Odczyt z macierzy sąsiedztwa zmusza algorytm do sprawdzenia każdej możliwej pary wierzchołków ($O(V^2)$), nawet w przypadku braku krawędzi. Zastosowanie listy następników skraca ten proces jedynie do przeglądu istniejących połączeń $O(V + E)$. Oczekuje się zatem przewagi listy na etapie inicjalizacji dla grafów rzadkich, jednak w ostatecznym rozrachunku dla grafów gęstych o czasie wykonania w obu wersjach zadecyduje sam koszt sortowania krawędzi w kolejce ($O(E \log E)$).
        \cite{Kruskal}
      \end{itemize}
    \subsubsection{Znalezienie najkrótszej ścieżki}
      Znalezienie najkrótszego pod względem sumy wag połączenia między dwoma wierzchołkami w grafie.
      \begin{itemize}
        \item{Algorytm Dijkstry}\\
Algorytm Dijkstry przeszukuje graf za pomocą strategii zachłannej, wyciągając z kolejki priorytetowej wierzchołki o najmniejszym koszcie dotarcia i relaksując ich sąsiadów. Podobnie jak w algorytmie Prima, reprezentacja ma kluczowy wpływ na etap weryfikacji sąsiedztwa. Użycie macierzy sąsiedztwa oznacza konieczność iterowania po wszystkich $V$ wierzchołkach dla każdego analizowanego węzła, co owocuje złożonością czasową w okolicach $O(V^2 \log V)$. Wersja na liście następników filtruje iteracje wyłącznie do fizycznych sąsiadów, osiągając optymalny czas $O(E \log V)$. Przewiduje się, że reprezentacja listowa zdecydowanie zdominuje macierzową w grafach rzadkich, natomiast w sieciach zbliżonych do grafów pełnych różnice w czasie wykonywania ulegną silnemu zatarciu.
\cite{Dijkstra}
        \item{Algorytm Forda-Bellmana}\\
Algorytm Forda-Bellmana, bazujący na programowaniu dynamicznym, opiera się na wykonaniu $V-1$ iteracji, w których próbuje zrelaksować wszystkie połączenia w grafie. Wybór reprezentacji determinuje tu drastyczne różnice w wydajności. W implementacji z macierzą sąsiedztwa mechanizm wymusza sprawdzanie w każdej z iteracji pełnej macierzy $V \times V$, co generuje fatalną w skutkach złożoność sześcienną $O(V^3)$. Zmiana struktury na listę następników sprawia, że w każdej głównej iteracji algorytm przebywa wyłącznie po istniejących krawędziach ($E$), redukując czas do poziomu $O(V \cdot E)$. Oczekuje się, że w przypadku tej metody reprezentacja listowa zmiażdży macierzową pod kątem wydajności na niemal każdym badanym zbiorze danych.
\cite{Ford-Bellman}
      \end{itemize}

 \subsection{Charakterystyka danych badawczych}
 Dane wykorzystane w badaniu zostały wygenerowane przy użyciu generatora liczb pseudolosowych. Podstawowym badanym typem danych dla wag krawędzi jest format liczb całkowitych (int, 4 bajty).
 Badanie obu algorytmów rozwiązujących ten sam problem zostało przeprowadzone na tych samych instancjach reprezentacji. Oba typy reprezentacji opisują ten sam badany graf.\\
 \\
 W badaniu uwzględniono trzy reprezentatywne poziomy gęstości grafu:
 \begin{itemize}
   \item{25\%}
   \item{50\%}
   \item{99\%}
 \end{itemize}
 Badania wykonano dla siedmiu reprezentatywnych rozmiarów grafu: 50, 100, 200, 400, 600, 800, 1000 wierzchołków.
 Dla każdej kombinacji parametrów badanie wykonano na 50 wygenerowanych instancjach problemu.
 \section{Warunki badawcze}
  \subsection{Środowisko badawcze}
  Badanie zostało przeprowadzone za pomocą autorskiego oprogramowania napisanego w języku
C++, skompilowanego za pomocą kompilatora Apple Clang w wersji 17.0.0 ,zgodnie ze standardem C++20, stosując standardową konfigurację. Specyfikacja platformy badawczej:
  \begin{itemize}
    \item{\textbf{Procesor:} Apple M4, 10 rdzeni (4 rdzenie wydajnościowe, 6 rdzeni energooszczędnych)}
    \item{\textbf{Pamięć operacyjna:} 16 GB zunifikowanej pamięci RAM}
  \end{itemize}
 \begin{thebibliography}{9}
  \bibitem{Prim}
 Thomas H. Cormen, Charles E. Leiserson, Ronald L Rivest, Clifford Stein: Wprowadzenie do algorytmów, PWN, 2024, s. 573-574.
  \bibitem{Kruskal}
  Thomas H. Cormen, Charles E. Leiserson, Ronald L Rivest, Clifford Stein: Wprowadzenie do algorytmów, PWN, 2024, rozdział 21.
  \bibitem{Dijkstra}
  Wikipedia Algorytm Dijkstry - dostęp 19.06.2026 (pl.wikipedia.org/wiki/Algorytm\_Dijkstry)
  \bibitem{Ford-Bellman}
 Thomas H. Cormen, Charles E. Leiserson, Ronald L Rivest, Clifford Stein: Wprowadzenie do algorytmów, PWN, 2024, s.664-665.
\end{thebibliography}
\end{document}
